\section*{LỜI NÓI ĐẦU}
\phantomsection\addcontentsline{toc}{section}{\numberline {}LỜI NÓI ĐẦU}

Đồ án này được thực hiện trong bối cảnh mạng 5G phi mặt đất (NTN) đang trở thành một trong những hướng phát triển trọng tâm của cộng đồng nghiên cứu viễn thông quốc tế. Kể từ khi 3GPP chính thức đưa hỗ trợ NTN vào Release 17 (năm 2022), câu hỏi về khả năng tương thích của các cơ chế giao thức vốn được thiết kế cho mạng mặt đất -- trong đó có HARQ -- với đặc thù độ trễ cao và liên kết LOS mạnh của kênh vệ tinh, trở thành một bài toán mở và có giá trị thực tiễn cao.

Công trình nghiên cứu của Tuninato và cộng sự (2025)~\cite{tuninato2025} là một trong những bài báo đầu tiên phân tích hệ thống sự đánh đổi giữa HARQ và RLC ARQ trong 5G NTN. Đồ án này được xây dựng dựa trên nền tảng đó, mở rộng thêm phân tích theo nhiều chiều: tổng hợp N$_{\min}$ cho bốn quỹ đạo × bốn cấu hình SCS, đánh giá hiệu quả năng lượng trên bit, và xác minh độc lập toàn bộ kết quả số bằng phân tích CDF Rician từ lý thuyết xác suất.

Em xin chân thành cảm ơn giảng viên hướng dẫn đã định hướng và hỗ trợ trong suốt quá trình thực hiện đồ án.

\vspace{12pt}
\begin{flushright}
\textit{Hà Nội, tháng 6 năm 2026}\\
\textit{Sinh viên thực hiện}\\
\vspace{24pt}
\textbf{Hoàng Tuấn Dương}
\end{flushright}
\cleardoublepage
